\subsection{Positive linear forms}

DEFINITION 2. --- Let A be an involutive algebra. A linear form \(\lambda\) on A is said to be positive if \(\lambda(a^{*}a) \in \mathbf{R}_{+}\) for all \(a \in \mathrm{A}\) .

If A is an involutive Banach algebra, we denote by \(A_{+}^{\prime}\) the set of continuous positive linear forms on A.

Let A be an involutive Banach algebra. The set \(A_{+}^{\prime}\) is a pointed convex cone in the real vector space of C-linear forms on A.

Lemma 2. --- Let A be an involutive Banach algebra and let \(\lambda\) be a positive linear form on A.

\begin{enumerate}
\def\labelenumi{\alph{enumi})}
\tightlist
\item
  For all \(a\) and \(b\) in A, we have \(\lambda(a^{*}b)=\overline{\lambda(b^{*}a)}\) and
\end{enumerate}

\[
\left| \lambda \left(a ^ {*} b\right) \right| ^ {2} \leqslant \lambda \left(a ^ {*} a\right) \lambda \left(b ^ {*} b\right);
\]

\begin{enumerate}
\def\labelenumi{\alph{enumi})}
\setcounter{enumi}{1}
\tightlist
\item
  If A is unital, then the linear form \(\lambda\) is continuous and its norm is equal to \(\lambda(1)\) .
\end{enumerate}

The map \((a,b)\mapsto\lambda(a^{*}b)\) is a positive Hermitian form on A; it therefore satisfies \(\lambda(a^{*}b)=\overline{\lambda(b^{*}a)}\) and \(|\lambda(a^{*}b)|^{2}\leqslant\lambda(a^{*}a)\lambda(b^{*}b)\) for all a and b in A (EVT, V, p. 2, remark and p. 3, prop. 2).

Let us prove \(b\) ). Let \(a \in A\) be a Hermitian element of norm \(< 1\) . The spectral radius of \(a\) is less than \(\|a\|\) , hence the spectrum of \(a\) is contained in the open unit disk D of C (th. 1 of I, p. 24). Let \(f\) be a holomorphic function on D such that \(f(z)^{2} = 1 - z\) for all \(z \in D\) (Lemma 1 of V, p. 435). Apply the holomorphic functional calculus to the element \(a\) and to the function \(f\) . The element \(b = f(a)\) satisfies \(b^{2} = 1 - a\) (I, p. 74, Thm. 5). By Prop. 3 of V, p. 435, we also have \(f(a)^{*} = f^{*}(a^{*}) = f(a)\) , hence \(b\) is Hermitian. It then follows that \(\lambda(1 - a) = \lambda(b^{*}b) \geqslant 0\) , whence \(\lambda(a) \leqslant \lambda(1)\) .

Now let \(b \in A\) be of norm \(< 1\) . The element \(b^*b\) is Hermitian and \(\|b^*b\| < 1\) , so applying \(a\) ) with \(a = 1\) , we find

\[
| \lambda (b) | ^ {2} \leqslant \lambda (1) \lambda (b ^ {*} b) \leqslant \lambda (1) ^ {2},
\]

with equality if \(b = 1\) . The linear form \(\lambda\) is therefore continuous, and its norm is equal to \(\lambda(1)\) . The assertion \(b\) ) is proved.

Example. --- Let X be a compact topological space and let A be the star algebra \(\mathcal{C}(\mathrm{X})\) . Positive linear forms on A are identified with positive measures on X (INT, III, p. 52, § 1, n° 6, th. 1).

Lemma 3. --- Let A be a unital involutive Banach algebra admitting an approximate unit (I, p. 120, def. 7). Let λ be a continuous positive linear form on A.

\begin{enumerate}
\def\labelenumi{\alph{enumi})}
\item
  For all a in A, we have \(\lambda(a^{*})=\overline{\lambda(a)}\) and \(|\lambda(a)|^{2}\leqslant\|\lambda\|\lambda(a^{*}a)\) ;
\item
  Let \(\widetilde{\mathrm{A}}\) be the involutive algebra obtained from A by adjoining a unit element, and let e be its unit element. There exists a continuous positive linear form \(\widetilde{\lambda}\) on \(\widetilde{\mathrm{A}}\) which extends \(\lambda\) and such that \(\widetilde{\lambda}(e) = \|\lambda\|\) ;
\item
  For all \(a\) and \(b\) in A, we have \(|\lambda(b^{*}ab)|\leqslant\|a\|\lambda(b^{*}b)\) .
\end{enumerate}

Let us prove \(a\) ). Let \(\mathfrak{F}\) be an approximate unit of A. Let \(a \in \mathrm{A}\) . Using Lemma 2, \(a\) ) and the definition of an approximate unit, one finds

\[
\lambda (a ^ {*}) = \lim _ {f, \mathfrak {F}} \lambda (f a ^ {*}) = \lim _ {f, \mathfrak {F}} \overline {{\lambda (a f ^ {*})}} = \lim _ {f, \mathfrak {F}} \overline {{\lambda ((f a ^ {*}) ^ {*})}} = \overline {{\lambda (a)}},
\]

whence (loc. cit.)

\[
| \lambda (a) | ^ {2} = \lim _ {f, \mathfrak {F}} | \lambda (f a) | ^ {2} \leqslant \lambda (a ^ {*} a) \operatorname * {limsup} _ {f, \mathfrak {F}} \lambda (f f ^ {*}) \leqslant \| \lambda \| \lambda (a ^ {*} a),
\]

since \(\mathfrak{F}\) is a filter on the unit ball of A.

To prove \(b\) ), one may assume that \(\lambda\) is not zero, and then that \(\| \lambda \| = 1\) . For \(a \in \mathrm{A}\) and \(z \in \mathbf{C}\) , set \(\widetilde{\lambda}(a + z \cdot e) = \lambda(a) + z\) . The map \(\widetilde{\lambda}\) is a continuous linear form on \(\widetilde{\mathrm{A}}\) which extends \(\lambda\) and satisfies \(\widetilde{\lambda}(e) = 1\) . It is positive: for all \(a \in \mathrm{A}\) and \(z \in \mathbf{C}\) , according to \(a\) ), one computes

\[
\begin{array}{r l} \widetilde {\lambda} ((a + z \cdot e) ^ {*} (a + z \cdot e)) & = \lambda (a ^ {*} a) + \overline {{z}} \lambda (a) + z \lambda (a ^ {*}) + | z | ^ {2} \\ & = | z + \lambda (a) | ^ {2} + \lambda (a ^ {*} a) - | \lambda (a) | ^ {2} \geqslant 0. \end{array}
\]

Finally, let \(b\) be an element of A. The map \(a \mapsto \widetilde{\lambda}(b^{*}ab)\) is a positive linear form on the unital involutive Banach algebra \(\widetilde{A}\) . It is therefore continuous with norm equal to its value at \(e\) (Lemma 2, \(b\) ), which is equal to \(\lambda(b^{*}b)\) , whence assertion \(c\) ).

PROPOSITION 4. — Let A be an involutive Banach algebra.

\begin{enumerate}
\def\labelenumi{\alph{enumi})}
\item
  For every Hilbert space E, every morphism of involutive algebras \(\varphi\colon A\to\mathcal{L}(E)\) and every vector \(x\in E\) , the linear form defined on A by \(\lambda(a)=\langle x|\varphi(a)x\rangle\) is a continuous positive linear form;
\item
  Let \(\lambda \in \mathrm{A}_{+}^{\prime}\) . The map \(f\) from \(\mathrm{A} \times \mathrm{A}\) to \(\mathbf{C}\) defined by \(f(a, b) = \lambda (a^{*}b)\) for all \((a, b) \in \mathrm{A} \times \mathrm{A}\) is a universally positive kernel on \(\mathrm{A}\) .
\end{enumerate}

Let us prove \(a\) ). For every \(a \in \mathrm{A}\) , we have

\[
\lambda (a ^ {*} a) = \langle x \mid \varphi (a ^ {*} a) x \rangle = \| \varphi (a) x \| ^ {2} \geqslant 0
\]

hence \(\lambda\) is a positive linear form on A; it is continuous since \(\varphi\) is continuous (Prop. 2 of I, p. 104).

Let us prove \(b\) ). The function \(f\) is continuous. For every \(n \in \mathbf{N}\) , every family \((a_i)_{0 \leqslant i \leqslant n}\) in A and every family \((t_i)_{0 \leqslant i \leqslant n}\) of complex numbers, we have

\[
\sum_ {i = 0} ^ {n} \sum_ {j = 0} ^ {n} \bar {t} _ {i} t _ {j} f (a _ {i}, a _ {j}) = \lambda \left(\left(\sum_ {i = 0} ^ {n} t _ {i} a _ {i}\right) ^ {*} \left(\sum_ {j = 0} ^ {n} t _ {j} a _ {j}\right)\right) \geqslant 0,
\]

whence the result (Def. 1 of V, p. 433).

DEFINITION 3. — Let A be an involutive Banach algebra. Let \(\lambda \in \mathrm{A}_{+}^{\prime}\) be a continuous positive linear form on A. A Hilbertian realization of \(\lambda\) is a triplet (E, \(x, \varphi\) ) consisting of a Hilbert space E, an element \(x\) of E, and a morphism of involutive algebras \(\varphi\) from A into \(\mathcal{L}(\mathrm{E})\) such that \(\lambda(b^{*}a) = \langle \varphi(b)x | \varphi(a)x \rangle\) for all \((a,b) \in \mathrm{A}^{2}\) .

If the set of elements \(\varphi(a)x\) for \(a\in\mathrm{A}\) is total in E, we say that (E, \(x,\varphi\) ) is a cyclic Hilbert-space realization of \(\lambda\) .

PROPOSITION 5. --- Let A be an involutive Banach algebra admitting an approximate unit. Every continuous positive linear form on A admits a cyclic Hilbert realization (E, x, φ). If A is unital, then there exists such a realization where the morphism φ is unital.

We may assume that A is a unital algebra (Lemma 3, b)).

Let \(\lambda\) be a positive continuous linear form on A. Let (E, \(g\) ) be a cyclic Hilbertian realization of the universally positive kernel \(f\) on A defined by \(f(a,b) = \lambda(a^{*}b)\) (prop. 4 and th. 1). The map \(g\) is continuous, its image is total in E, and we have \(\lambda(a^{*}b) = \langle g(a)|g(b)\rangle\) for all \((a,b) \in \mathrm{A}^2\) . Set \(x = g(1) \in \mathrm{E}\) .

The map g from A into E is linear: indeed, for \((a,b,c)\in\mathrm{A}^{3}\) and \((s,t)\in\mathbf{C}^{2}\) , it follows that

\[
\begin{array}{c} \langle g (c) | g (s a + t b) \rangle = \lambda (c ^ {*} (s a + t b)) \\ = s \lambda (c ^ {*} a) + t \lambda (c ^ {*} b) = \langle g (c) | s g (a) + t g (b) \rangle , \end{array}
\]

whence the assertion since the image of \(g\) is total in E.

In particular, the image F of g is a dense vector subspace of E. The kernel of g is a left ideal of A: if \(g(b) = 0\) , then for all a and c in A, it follows that

\[
\langle g (a b) \mid g (c) \rangle = \lambda ((a b) ^ {*} c) = \lambda (b ^ {*} (a ^ {*} c)) = \langle g (b) \mid g (a ^ {*} c) \rangle = 0,
\]

hence \(g(ab) = 0\) since F is dense in E.

Let \(a \in \mathrm{A}\) . Since the kernel of \(g\) is a left ideal of \(\mathrm{A}\) , there exists a linear map \(\widetilde{\varphi}(a) \colon \mathrm{F} \to \mathrm{F}\) such that \(\widetilde{\varphi}(a)(g(b)) = g(ab)\) for all \(b \in \mathrm{A}\) . Moreover, for every \(b \in \mathrm{A}\) , it follows that

\[
\| \widetilde {\varphi} (a) (g (b)) \| ^ {2} = \| g (a b) \| ^ {2} = \lambda \left(b ^ {*} a ^ {*} a b\right) \leqslant \| a ^ {*} a \| \lambda \left(b ^ {*} b\right) \leqslant \| a \| ^ {2} \| g (b) \| ^ {2}
\]

(lemma 3, c)), hence \(\widetilde{\varphi}(a)\) is continuous and of norm \(\leqslant \|a\|\) . Consequently, there exists a unique endomorphism \(\varphi(a) \in \mathcal{L}(\mathrm{E})\) which induces \(\widetilde{\varphi}(a)\) by passage to the subspaces.

Let \(a\) and \(b\) be in A. We have

\[
(\varphi (a) \circ \varphi (b)) (g (c)) = g (a b c) = \varphi (a b) (g (c))
\]

for all c in A, hence \(\varphi(ab) = \varphi(a) \circ \varphi(b)\) since F is dense in E.

The map \(\varphi\) from A into \(\mathcal{L}(\mathrm{E})\) is linear: indeed, for every \((a,b)\in\mathrm{A}^{2}\) and every \((s,t)\in\mathbf{C}^{2}\) , we have

\[
\varphi (s a + t b) (g (c)) = g ((s a + t b) c) = (s \varphi (a) + t \varphi (b)) g (c)
\]

for all \(c \in A\) , whence \(\varphi(sa + tb) = s\varphi(a) + t\varphi(b)\) since F is dense in E. Since \(\|\varphi(b)\| \leqslant \|b\|\) for all \(b \in A\) , the map \(\varphi\) is continuous. We also have \(\varphi(1) = 1_{\mathrm{E}}\) since \(\varphi(1)(g(a)) = g(a)\) for all \(a \in A\) and since F is dense in E.

Finally, let a, b, c be in A. We have

\[
\begin{array}{c} \langle g (c) \mid \varphi (a ^ {*}) (g (b)) \rangle = \langle g (c) \mid g (a ^ {*} b) \rangle = \lambda (c ^ {*} a ^ {*} b) \\ = \langle g (a c) \mid g (b) \rangle = \langle \varphi (a) (g (c)) \mid g (b) \rangle \end{array}
\]

whence \(\varphi(a^{*})=\varphi(a)^{*}\) since F is dense in E.

In conclusion, the map \(\varphi\) is a continuous morphism of involutive algebras from A into \(\mathcal{L}(\mathrm{E})\) and \(g(a) = \varphi(a)x\) for all \(a \in \mathrm{A}\) ; consequently, the triple \((\mathrm{E}, x, \varphi)\) is a cyclic Hilbert realization of \(\lambda\) .

PROPOSITION 6. --- Let A be an involutive Banach algebra and let λ be a continuous positive linear form on A. Let (E₁, x₁, φ₁) and (E₂, x₂, φ₂) be Hilbertian realizations of λ. Suppose that (E₁, x₁, φ₁) is a cyclic Hilbertian realization.

\begin{enumerate}
\def\labelenumi{\alph{enumi})}
\item
  There exists a unique continuous linear map u from \(E_{1}\) to \(E_{2}\) which is a morphism of representations from \(\varphi_{1}\) into \(\varphi_{2}\) (I, p. 11) and which satisfies \(u(\varphi_{1}(a)x_{1}) = \varphi_{2}(a)x_{2}\) for all \(a \in A\) ;
\item
  The linear map u is isometric;
\item
  If \(\left(\mathrm{E}_{2},x_{2},\varphi_{2}\right)\) is cyclic, then u is an isomorphism;
\item
  If \(\left(\mathrm{E}_2,x_2,\varphi_2\right)\) is cyclic and if A is unital, then \(u(x_{1}) = x_{2}\) .
\end{enumerate}

For \(j \in \{1,2\}\) , define \(\gamma_j \colon A \to E_j\) by \(\gamma_j(a) = \varphi_j(a)x_j\) for all \(a \in A\) . By definition, the pairs \((E_j, \gamma_j)\) are Hilbert realizations of the universally positive kernel \(f\) on \(A\) defined by \((a,b) \mapsto \lambda(a^*b)\) . The Hilbert realization \((E_1, \gamma_1)\) is cyclic; by prop. 1 of V, p. 434, there therefore exists a unique continuous linear map \(u \colon E_1 \to E_2\) such that \(\gamma_2 = u \circ \gamma_1\) , and this map is isometric. Moreover, if \((E_2, x_2, \varphi_2)\) is also cyclic, then \(u\) is an isomorphism. To prove \(a)\) , \(b)\) and \(c)\) , it suffices to prove that \(u\) is a morphism of representations from \(\varphi_1\) into \(\varphi_2\) .

Let \(a \in \mathrm{A}\) . For every \(b \in \mathrm{A}\) , we have

\[
\begin{array}{c} (u \circ \varphi_ {1} (a)) (\gamma_ {1} (b)) = (u \circ \varphi_ {1} (a b)) x _ {1} = (u \circ \gamma_ {1}) (a b) \\ = \gamma_ {2} (a b) = \varphi_ {2} (a) (\gamma_ {2} (b)) = \varphi_ {2} (a) (u (\gamma_ {1} (b))), \end{array}
\]

therefore the continuous linear maps \(u \circ \varphi_1(a)\) and \(\varphi_2(a) \circ u\) coincide on the subspace of \(E_1\) generated by the image of \(\gamma_1\) ; this subspace is dense in \(E_1\) by hypothesis, hence \(u \circ \varphi_1(a) = \varphi_2(a) \circ u\) .

Let us finally prove \(d\) ). Suppose therefore that A is unital and that \((\mathrm{E}_2, x_2, \varphi_2)\) is cyclic. There exists a cyclic Hilbert-space realization \((\mathrm{E}_3, x_3, \varphi_3)\) of \(\lambda\) such that \(\varphi_3\) is a unital morphism (prop. 5). Let \(j \in \{1, 2\}\) . Applying the preceding to \((\mathrm{E}_j, x_j, \varphi_j)\) and \((\mathrm{E}_3, x_3, \varphi_3)\) , one sees that there exists an isometric isomorphism \(\widetilde{u}_j\) from \(\mathrm{E}_j\) into \(\mathrm{E}_3\) such that \(\widetilde{u}_j \circ \varphi_j(a) = \varphi_3(a) \circ \widetilde{u}_j\) for all \(a \in \mathrm{A}\) . Taking \(a = 1_{\mathrm{A}}\) , one sees that \(\varphi_j\) is unital. One then has \(x_2 = \varphi_2(1_{\mathrm{A}})x_2 = u(\varphi_1(1_{\mathrm{A}})x_1) = u(x_1)\) .

Remark. --- Keep the notations of the proposition. It is possible that \(u(x_{1})\) is different from \(x_{2}\) (exercise 3, b) of V, p. 493). However, the triple \((\mathrm{E}_{2}, u(x_{1}), \varphi_{2})\) is also a Hilbertian realization of \(\lambda\) .

